%%%%%%%%%%%%%%%%%%%%%%%%%%%%%%%%%%%%%%%%%%%%%%%%%%%%%%%%%%%%%%%%%%%%%%
% Lớp 10 | Phương pháp tọa độ trong mặt phẳng (𝑂𝑥𝑦) | Tích vô hướng (theo tọa độ)
% Dạng 0H9?2-4: Điều kiện vuông góc
% Số câu: 5  (sinh tự động bởi tach_cau.py; từ nay thầy sửa trực tiếp file này)
%%%%%%%%%%%%%%%%%%%%%%%%%%%%%%%%%%%%%%%%%%%%%%%%%%%%%%%%%%%%%%%%%%%%%%

%%%%%%%%%%%%%%%% TRẮC NGHIỆM 4 PHƯƠNG ÁN %%%%%%%%%%%%%%%%

% nguồn: MaHoaNH_TracNghiem_Lop10
\begin{ex}%[0H9H2-4]%[ma:64ff97e3]
	Trong mặt phẳng tọa độ $Oxy$, cho hai vectơ $\overrightarrow{u}=\left(1;6\right)$ và $\overrightarrow{v}=\left(m;-3\right)$. Giá  trị của $m$ để $\overrightarrow{u}$ vuông góc với $\overrightarrow{v}$ là
	\choice 
		{$m=-18$} 
		{$m=\dfrac{1}{2}$} 
		{$m=-\dfrac{1}{2}$} 
		 {\True $m=18$}
	 \loigiai{
		$\overrightarrow{u}\perp  \overrightarrow{v}\Leftrightarrow \overrightarrow{u}\cdot \overrightarrow{v}=0 \Leftrightarrow 1\cdot m +6\cdot \left(-3\right)=0\Leftrightarrow m= 18$.
	} 
\end{ex}


%%%%%%%%%%%%%%%% ĐÚNG/SAI %%%%%%%%%%%%%%%%

% nguồn: MaHoaNH_DungSai_Lop10
\begin{ex}%[0H9H2-4]%[ma:816d8658]
	Trong hệ trục tọa độ $Oxy$, cho ba véc-tơ $\overrightarrow{a}=(3;-1)$, $\overrightarrow{b}=(1;-2)$,  $\overrightarrow{c}=(-1;7)$. Xét tính đúng - sai của các khẳng định bên dưới.
	 \choiceTF
		 {\True $\overrightarrow{u}=\overrightarrow{a}-2\overrightarrow{b}+3\overrightarrow{c}=(-2;24)$}
		{\True $\left|\overrightarrow{a}\right|^2=2\left|\overrightarrow{b}\right|^2$} 
		 {$\overrightarrow{a}\cdot\left(\overrightarrow{b}+\overrightarrow{c}\right)=12$}
		{\True Véc-tơ $\overrightarrow{x}=\overrightarrow{b}+5\overrightarrow{i}$  vuông góc với véc-tơ $\overrightarrow{a}$}
	\loigiai{\begin{itemchoice}
		\itemch  Giả sử $\overrightarrow{u}=(x_0;y_0;z_0)$. Khi đó
		$\heva{&x_0=3-2\cdot 1+ 3\cdot (-1)=-2\\&y_0=-1-2\cdot (-2)+3\cdot7=24}$.  Vậy $\overrightarrow{u}=( -2;24)$ là đúng.
		\itemch Ta có  $\left|\overrightarrow{a}\right|=\sqrt{10}$, $\left|\overrightarrow{b}\right|=\sqrt{5}$ nên $\left|\overrightarrow{a}\right|^2=2\left|\overrightarrow{b}\right|^2$ là đúng.
		\itemch Ta có  $\overrightarrow{b}+\overrightarrow{c}=(0;5)$ nên $\overrightarrow{a}\cdot\left(\overrightarrow{b}+\overrightarrow{c}\right)=-5$, suy ra $\overrightarrow{a}\cdot\left(\overrightarrow{b}+\overrightarrow{c}\right)=12$ là sai.
		\itemch Ta có $\overrightarrow{x}=\overrightarrow{b}+5\overrightarrow{i}=(1;3)$, suy ra $\overrightarrow{x}\cdot\overrightarrow{a}=0$ nên véc-tơ $\overrightarrow{x}=\overrightarrow{b}+5\overrightarrow{i}$ vuông góc với  véc-tơ $\overrightarrow{a}$ là đúng.
\end{itemchoice}}
\end{ex}

% nguồn: MaHoaNH_DungSai_Lop10
\begin{ex}%[0H9H2-4]%[ma:35879115]
	Trong hệ tọa độ $Oxy$, các khẳng định sau  đây đúng hay sai?
	 \choiceTF
		{\True Hai véc-tơ $\overrightarrow{i}$, $\overrightarrow{j}$ luôn vuông  góc với nhau}
		{\True Cặp  véc-tơ $\overrightarrow{a}=(-4;-6)$ và $\overrightarrow{b}=(-3;2)$ vuông góc với nhau}
		{Hai véc-tơ $\overrightarrow{u}=\left(\dfrac{1}{2};-5 \right)$, $\overrightarrow{v}=\left(k;-4  \right)$ vuông góc với nhau khi $k=-4$}
		{Biết rằng $\overrightarrow{u}=(4;1)$, $\overrightarrow{v}=(1;4)$  và $\overrightarrow{a}=\overrightarrow{u}+m\cdot\overrightarrow{v}$ với $m\in\mathbb{R}$. Khi đó $m=-3$ thì $\overrightarrow{a}$ vuông góc với $\overrightarrow{b}=(1;-1)$}
	 \loigiai{
		\begin{itemchoice}
			\itemch Hai véc-tơ $\overrightarrow{i}$, $\overrightarrow{j}$ là các véc-tơ đơn vị lần lượt nằm trên trục $Ox$, $Oy$ nên chúng luôn vuông góc  với nhau.
			\itemch Ta có $\overrightarrow{a}\cdot\overrightarrow{b}=0$  nên $\overrightarrow{a}\perp \overrightarrow{b}$.
			\itemch Ta có $\overrightarrow{u}\cdot\overrightarrow{v}=0 \Leftrightarrow \dfrac{1}{2}\cdot k+(-5)\cdot  (-4)=0\Leftrightarrow k=-40$.
			\itemch Ta có  $\overrightarrow{a}=\overrightarrow{u}+m\cdot\overrightarrow{v}=(4+m; 1+4m)$.\\
			Khi đó $\overrightarrow{a}\cdot\overrightarrow{b}=0 \Leftrightarrow 4+m-1-4m=0 \Leftrightarrow  m=1$.
		\end{itemchoice}
	 }
\end{ex}

% nguồn: MaHoaNH_DungSai_Lop10
\begin{ex}%[0H9V2-4]%[ma:e8907a34]
	Trong hệ tọa độ $Oxy$, cho $\triangle ABC$ có $A(2;4)$, $B(2;-2)$, $C(-4;1)$ và $\overrightarrow{v}=(m^2+1;2-m)$. Gọi $K(a;b)$ là hình chiếu vuông  góc của điểm $A$ trên cạnh $BC$. Trong các khẳng định sau, khẳng định nào đúng, khẳng định nào sai?
	 \choiceTF
		{\True $\overrightarrow{v}\perp\overrightarrow{AB}$ khi $m=2$} 
		{Có $2$ điểm $M$ thuộc trục hoành sao  cho $\triangle MBC$ vuông tại $B$}
		{\True Trực tâm của $\triangle ABC$ là $H\left(\dfrac{1}{2};  1\right)$}
		{$3a+b=2$} 
	\loigiai{ 
		\begin{itemchoice}
			\itemch Ta  có $\overrightarrow{AB}=(0;-6)$.\\
			Khi đó $\overrightarrow{v}\cdot\overrightarrow{AB}=0 \Leftrightarrow  0\cdot (m^2+1)-6\cdot (2-m)=0 \Leftrightarrow m=2$.
			 \itemch Gọi $M(x_M;0)\in Ox$.\\
			Ta có $\overrightarrow{BM}=(x_M-2;2)$,  $\overrightarrow{BC}=(-6;3)$.\\
			Khi đó $\triangle MBC$ vuông tại  $B$
			\allowdisplaybreaks 
			\begin{eqnarray*}
				&\Leftrightarrow & \overrightarrow{BM}\cdot\overrightarrow{BC}=0  \Leftrightarrow -6x_M+18=0\\
				&\Leftrightarrow  & x_M=3 \Rightarrow M(3;0)
			\end{eqnarray*}
			Vậy có  $1$ điểm $M(3;0)$ thuộc trục hoành sao cho $\triangle MBC$ vuông tại $B$.
			 \itemch Ta có $\overrightarrow{BC}=(-6;3)$, $\overrightarrow{AB}=(0;-6)$.\\
			Giả sử tọa  độ trực tâm $H$ của $\triangle ABC$ là $H(x;y)$, ta có
			 \allowdisplaybreaks
			\begin{eqnarray*}
				& & \heva{& AH\perp BC\\& CH\perp AB} \Leftrightarrow\heva{& \overrightarrow{AH}\cdot\overrightarrow{BC}=0\\&  \overrightarrow{CH}\cdot\overrightarrow{AB}=0}\\
				&\Leftrightarrow & \heva{& -6(x-2)+3(y-4)=0\\& 0(x+4)-6(y-1)=0} \Leftrightarrow\heva{& x=\dfrac{1}{2}\\& y=1.} 
			\end{eqnarray*}
			Vậy $H\left(\dfrac{1}{2};  1\right)$.
			\itemch Ta có $\overrightarrow{AK}=(a-2; b-4)$, $\overrightarrow{BC}=(-6;3)$, $\overrightarrow{BK}=(a-2; b+2)$.\\ 
			Theo giả  thiết, ta có
			\allowdisplaybreaks 
			\begin{eqnarray*}
				 & & \heva{& \overrightarrow{AK}\cdot\overrightarrow{BC}=0\\& \exists k\in [0;1] \colon\overrightarrow{BK}=k\cdot\overrightarrow{BC}} \Leftrightarrow \heva{& -6a+3b=0\\& a-2=-6k\\& b+2=3k\\& k\in [0;1]}\\
				&\Leftrightarrow & \heva{& -6a+3b=0\\& a+6k=2\\& b-3k=-2\\& k\in [0;1]} \Leftrightarrow \heva{& a=-\dfrac{2}{5}\\& b=-\dfrac{4}{5}\\&  k=\dfrac{2}{5} \text{ (thỏa mãn)}.}
			\end{eqnarray*}
			Vậy  $3a+b=-2$.
		\end{itemchoice}
	} 
\end{ex}


%%%%%%%%%%%%%%%% TỰ LUẬN %%%%%%%%%%%%%%%%

% nguồn: MaHoaNH_TuLuan_Lop10
\begin{bt}%[0H9Y2-4]%[0H9Y2-5]%[ma:0e8eb5bf]
	Trong mặt phẳng tọa độ $Oxy$, cho điểm $A(3;-1)$, đường thẳng $\Delta_1\colon x-3y+15=0$ và đường  thẳng $\Delta_2\colon 2x-y-100=0$.
	\begin{enumerate}
		\item Tính góc giữa đường thẳng  $\Delta_1$ và $\Delta_2$.
		\item  Tính khoảng cách từ điểm $A$ đến đường thẳng $\Delta_2$.
	\end{enumerate}
	\loigiai{ 
		\begin{enumerate}
			\item $\Delta_1$ có véc-tơ pháp tuyến  $\overrightarrow{n}_1=(1;-3)$ và $\Delta_2$ có véc-tơ pháp tuyến $\overrightarrow{n}_2=(2;-1)$.\\
			 $\cos(\Delta_1;\Delta_2)=\left|\cos(\overrightarrow{n}_1;\overrightarrow{n}_2)\right|=\dfrac{|1\cdot2+(-3)\cdot(-1)|}{\sqrt{1^2+(-3)^2}\cdot\sqrt{2^2+(-1)^2}}=\dfrac{\sqrt{2}}{2}$.\\
			Suy ra $(\Delta_1;\Delta_2)=45^\circ$. 
			\item  $\mathrm{d}(A,\Delta_2)=\dfrac{|2\cdot3+1-100|}{\sqrt{2^2+(-1)^2}}=\dfrac{93\sqrt{5}}{5}$.
		\end{enumerate}
	} 
\end{bt}
